\section{Введение}
\label{sec:Chapter0} \index{Chapter0}

\sethlcolor{yellow}

Разработка цифровой интегральной схемы --- это многоэтапный процесс, на каждой стадии
которого требуется верификация, то есть проверка соответствия полученного описания техническому
заданию или результату предыдущего этапа. Одним из ранних и наиболее важных этапов является
построение модели схемы на уровне регистровых передач (англ. Register Transfer Level, RTL),
которая описывает поведение схемы в терминах элементов памяти (регистров) и передачи данных между ними.
Для описания RTL-моделей используются языки описания аппаратуры
(англ. Hardware Description Languages, HDL), среди которых наибольшее распространение
получили Verilog и его расширение SystemVerilog [1].

Одним из самых распространённых шаблонов проектирования в HDL является конечный автомат
(англ. Finite State Machine, FSM). Конечный автомат позволяет сохранять текущее состояние
исполнения в специальном регистре и обновлять его на последующих тактах сигнала синхронизации.
Конечные автоматы дают возможность увеличить тактовую частоту схемы за счёт разбиения
вычисления на несколько тактов, а также реализовать ожидание внешних событий в течение
произвольного времени. Практически любое нетривиальное управляющее устройство цифровой схемы
содержит один или несколько конечных автоматов.

Несмотря на концептуальную простоту модели, реализация конечных автоматов на практике
сопряжена с риском возникновения ошибок, которые трудно обнаружить традиционными методами
верификации, основанными на моделировании. К таким ошибкам относятся недостижимые состояния
(в которые автомат не может попасть ни при каком входном воздействии),
мёртвые переходы (условие которых никогда не выполняется), дедлоки (блокировки) и лайвлоки 
(<<живые>> блокировки). Покрытие всех подобных сценариев стандартным симуляционным тестированием
в общем случае нетривиально, поскольку оно требует написания соответствующего тестового окружения.

Формальные методы верификации свойств (англ. Formal Property \\ Verification, FPV)
в принципе позволяют обнаруживать такие проблемы, однако требуют от инженеров~-~верификаторов
ручного написания свойств и значительных трудозатрат, что ограничивает их применимость
на ранних стадиях проектирования. Методы статического анализа конечных автоматов
дополняют FPV, обеспечивая автоматическое обнаружение структурных аномалий в графе автомата.
Это делает их особенно удобными для проверок на ранних этапах проектирования и для
интеграции в процессы непрерывной интеграции.

Данная работа посвящена разработке инструмента SVAN-FSM для автоматического обнаружения и
анализа конечных автоматов в описаниях на языке SystemVerilog. Инструмент разработан
в рамках проекта SVAN [3] — системы статического анализа SystemVerilog-описаний.
В основе предложенного подхода лежит извлечение конечного автомата из списка
логических элементов (netlist), получаемого после обработки исходного кода компилятором,
построение графового представления автомата и применение к нему алгоритмов статического
анализа для поиска типичных ошибок. Особое внимание уделяется сопоставлению результатов
анализа с исходным кодом, что существенно улучшает интерпретируемость полученных
данных для пользователя.
